\documentclass[10pt,letterpaper,oneside]{report}
\usepackage[T1]{fontenc}
\usepackage[utf8]{inputenc}
\usepackage{lmodern}
\usepackage{microtype}
\usepackage[margin=0.85in,headheight=15pt]{geometry}
\usepackage{enumitem,booktabs,longtable,tabularx,array,makecell}
\usepackage{fancyhdr,lastpage,xcolor,hyperref,titlesec}
\usepackage{listings}
\usepackage[most]{tcolorbox}
\definecolor{ink}{HTML}{182D40}
\definecolor{accent}{HTML}{126B77}
\definecolor{pale}{HTML}{EEF5F5}
\definecolor{grayblue}{HTML}{677887}
\definecolor{codebg}{HTML}{F5F7F8}
\hypersetup{colorlinks=true,linkcolor=accent,urlcolor=accent,pdfauthor={C23 Engineering Companion},pdfsubject={ISO/IEC 9899:2024 educational engineering companion}}
\titleformat{\chapter}[display]{\normalfont\huge\bfseries\color{ink}}{\filleft\large\color{accent}CHAPTER \thechapter}{1ex}{\titlerule[1pt]\vspace{1ex}}
\titleformat{\section}{\Large\bfseries\color{ink}}{\thesection}{0.7em}{}
\titleformat{\subsection}{\large\bfseries\color{accent}}{\thesubsection}{0.65em}{}
\setlist[itemize]{leftmargin=1.3em,itemsep=2pt,topsep=3pt}
\setlist[enumerate]{leftmargin=1.6em,itemsep=3pt,topsep=3pt}
\setlength{\parskip}{5pt}
\setlength{\parindent}{0pt}
\renewcommand{\arraystretch}{1.18}
\pagestyle{fancy}\fancyhf{}\fancyhead[L]{\small\sffamily C23 ENGINEERING SERIES}\fancyhead[R]{\small\sffamily\color{accent}\DocShort}\fancyfoot[L]{\small\color{grayblue}ISO/IEC 9899:2024 companion}\fancyfoot[R]{\small\thepage\ / \pageref{LastPage}}
\newtcolorbox{keybox}{breakable,colback=pale,colframe=accent,boxrule=.6pt,arc=1mm,left=8pt,right=8pt,top=6pt,bottom=6pt}
\newtcolorbox{warnbox}{breakable,colback=codebg,colframe=grayblue,boxrule=.6pt,arc=1mm,left=8pt,right=8pt,top=6pt,bottom=6pt}
\lstdefinestyle{cstyle}{language=C,basicstyle=\ttfamily\footnotesize,keywordstyle=\color{accent}\bfseries,commentstyle=\color{grayblue}\itshape,stringstyle=\color{ink},backgroundcolor=\color{codebg},frame=single,rulecolor=\color{grayblue!40},breaklines=true,showstringspaces=false,tabsize=4,columns=fullflexible,keepspaces=true,upquote=true}
\lstset{style=cstyle}
\newcommand{\refcl}[1]{\textsf{\footnotesize[ISO/IEC 9899:2024, #1]}}
\newcommand{\standard}{\textbf{[STANDARD]}}
\newcommand{\optional}{\textbf{[OPTIONAL]}}
\newcommand{\guidance}{\textbf{[GUIDANCE]}}
\newcommand{\testtag}{\textbf{[TEST]}}
\newcommand{\term}[1]{\textbf{#1}}
\newcommand{\DocShort}{IMPLEMENTATION}
\begin{document}
\begin{titlepage}
\color{ink}\vspace*{1.2in}
{\sffamily\small\color{accent}\bfseries TECHNICAL DOCUMENTATION SUITE \hfill EDITION 2026-10}\par
\vspace{.3in}\rule{\linewidth}{1.5pt}\par\vspace{.22in}
{\LARGE\bfseries C23 Hands-On Implementation Handbook}\par\vspace{.22in}
{\Large\color{grayblue}Progressive labs, test programs and capstone deliverables}\par\vspace{.4in}
\begin{keybox}\textbf{Standards basis:} ISO/IEC 9899:2024 (C23). Independent educational and engineering companion. Not an official ISO interpretation, full specification, certification or conformance assessment.\end{keybox}
\vfill{\large\bfseries C23 Software Engineering and Language Standards}\par\vspace{.15in}{\color{grayblue}Prepared October 2026 \quad | \quad Source edition: October 2024}\par
\end{titlepage}
\pagenumbering{roman}\tableofcontents\clearpage\pagenumbering{arabic}

\chapter*{Sources and interpretation key}
\addcontentsline{toc}{chapter}{Sources and interpretation key}
\textbf{Primary source:} ISO/IEC 9899:2024(en), \emph{Information technology --- Programming languages --- C}, fifth edition (2024). Clause and annex identifiers refer to the supplied licensed standard. This handbook contains original explanations, designs and examples; it does not reproduce the standard. The licensed standard remains the authoritative source for exact wording and conformance decisions.\par
\textbf{Secondary engineering sources:} \url{https://gcc.gnu.org/onlinedocs/gcc/Standards.html}; \url{https://clang.llvm.org/c_status.html}; \url{https://cmake.org/cmake/help/latest/prop_tgt/C_STANDARD.html}. Compiler support and diagnostics are release-specific: verify the installed compiler, target and library before relying on a feature.\par
\begin{keybox}\textbf{Interpretation categories.} \textbf{[STANDARD]} summarizes a requirement, constraint, semantic rule or defined facility and cites the applicable clause. \textbf{[OPTIONAL]} describes conditional or implementation-dependent support. \textbf{[GUIDANCE]} records an engineering recommendation, not an ISO mandate. \textbf{[TEST]} is an original check or exercise; a passing test is evidence but is not proof of full conformity.\end{keybox}
\textbf{Copyright boundary:} The supplied PDF is individually licensed. Do not redistribute it with this suite; distribute only these independently written companion materials in accordance with applicable rights.\par


\chapter{Reproducible workbench}
\section{Tooling contract}
Use a C23-aware compiler, a compatible C library and a build system that rejects dialect downgrade. The minimal practical stack contains GCC or Clang, CMake, a unit-test runner, a debugger, static analysis, and an address/undefined-behavior sanitizer where supported. Compiler support may be incomplete; record version-specific limitations instead of silently changing test expectations.
\begin{lstlisting}[language=bash]
cc --version
cc -std=c23 -Wall -Wextra -Wpedantic -Wconversion -c main.c
# Enable sanitizers only where available for the actual target.
cc -std=c23 -g -O1 -fsanitize=address,undefined main.c -o lab
./lab
\end{lstlisting}
\section{Repository layout}
\begin{lstlisting}[language=]
c23-workshop/
  CMakeLists.txt
  include/                   # API contracts
  src/                       # implementations
  tests/                     # positive and negative paths
  probes/                    # isolated C23 feature tests
  fixtures/                  # binary inputs; source-controlled
  docs/                      # behavior register, evidence
  .github/workflows/         # optional CI
\end{lstlisting}
\chapter{Reference workshop application: bounded record decoder}
\section{Protocol specification [GUIDANCE]}
Implement a wire format: one byte \texttt{kind}, two bytes \texttt{length} (network/big-endian), followed by exactly \texttt{length} payload bytes. The decoder rejects extra/short frames, does not assume host byte order, and borrows the input buffer (no allocation or ownership transfer).
\begin{lstlisting}
#include <stdint.h>
#include <stddef.h>
#include <stdbool.h>
struct Record {
    uint8_t kind;
    const uint8_t *payload;
    size_t length;
};
bool decode_record(const uint8_t *src, size_t src_size,
                   struct Record *out) {
    if (src == NULL || out == NULL || src_size < 3) return false;
    size_t length = ((size_t)src[1] << 8) | (size_t)src[2];
    if (length != src_size - 3) return false;
    struct Record next = { .kind = src[0],
                           .payload = src + 3, .length = length };
    *out = next;
    return true;
}
\end{lstlisting}
\guidance\ The output is a borrowed view; it becomes invalid when the source buffer's lifetime ends. If a zero-length payload is valid, \texttt{payload} may point one past the fixed-size header; do not dereference it when \texttt{length == 0}.
\section{Reference unit tests}
\begin{lstlisting}
#include <assert.h>
static void test_decoder(void) {
    const uint8_t ok[] = { 7, 0, 2, 0xAB, 0xCD };
    struct Record r = {0};
    assert(decode_record(ok, sizeof ok, &r));
    assert(r.kind == 7 && r.length == 2);
    assert(r.payload[0] == 0xAB && r.payload[1] == 0xCD);
    const uint8_t short_frame[] = { 7, 0, 4, 0xAB };
    assert(!decode_record(short_frame, sizeof short_frame, &r));
}
\end{lstlisting}
\chapter{Progressive laboratory catalog}
\section{Lab 1: Orientation and dialect}
\textbf{Standards map:} \refcl{1--4, 5.2}.\\
\textbf{Learning goal:} Identify compiler, libc, ABI and \_\_STDC\_VERSION\_\_.\\
\textbf{Procedure:} Compile a hello-world program in C17 and C23; retain version logs.\\
\textbf{Acceptance evidence:} Record differences and the selected conformance mode.
\begin{enumerate}
\item Create a clean branch and save the compiler/runtime version and invocation.
\item Build one success case and one intentionally failing test case.
\item Record expected diagnostic or runtime behavior without asserting unsupported guarantees.
\item Deliver source, commands, test output, and a one-paragraph standard-versus-policy note.
\end{enumerate}

\section{Lab 2: Lexing and preprocessor}
\textbf{Standards map:} \refcl{5.2.1, 6.4, 6.10}.\\
\textbf{Learning goal:} Understand expansion versus interpretation.\\
\textbf{Procedure:} Inspect -E output of nested macros and conditional inclusion.\\
\textbf{Acceptance evidence:} Explain every injected token.
\begin{enumerate}
\item Create a clean branch and save the compiler/runtime version and invocation.
\item Build one success case and one intentionally failing test case.
\item Record expected diagnostic or runtime behavior without asserting unsupported guarantees.
\item Deliver source, commands, test output, and a one-paragraph standard-versus-policy note.
\end{enumerate}

\section{Lab 3: Basic types and promotions}
\textbf{Standards map:} \refcl{6.2.5, 6.3}.\\
\textbf{Learning goal:} Trace integer ranks and conversions.\\
\textbf{Procedure:} Print sizeof and limits; write mixed signed/unsigned assertions.\\
\textbf{Acceptance evidence:} Expected and actual values match at boundary cases.
\begin{enumerate}
\item Create a clean branch and save the compiler/runtime version and invocation.
\item Build one success case and one intentionally failing test case.
\item Record expected diagnostic or runtime behavior without asserting unsupported guarantees.
\item Deliver source, commands, test output, and a one-paragraph standard-versus-policy note.
\end{enumerate}

\section{Lab 4: Lifetime and storage}
\textbf{Standards map:} \refcl{6.2.4}.\\
\textbf{Learning goal:} Prevent invalid pointers.\\
\textbf{Procedure:} Diagram stack, static and allocated object lifetimes.\\
\textbf{Acceptance evidence:} No pointers consumed after lifetime ends.
\begin{enumerate}
\item Create a clean branch and save the compiler/runtime version and invocation.
\item Build one success case and one intentionally failing test case.
\item Record expected diagnostic or runtime behavior without asserting unsupported guarantees.
\item Deliver source, commands, test output, and a one-paragraph standard-versus-policy note.
\end{enumerate}

\section{Lab 5: Arrays and pointer arithmetic}
\textbf{Standards map:} \refcl{6.5.7}.\\
\textbf{Learning goal:} Prove legal element access.\\
\textbf{Procedure:} Implement indexed byte slice with checked offset/count.\\
\textbf{Acceptance evidence:} Reject offset > length and overlong reads.
\begin{enumerate}
\item Create a clean branch and save the compiler/runtime version and invocation.
\item Build one success case and one intentionally failing test case.
\item Record expected diagnostic or runtime behavior without asserting unsupported guarantees.
\item Deliver source, commands, test output, and a one-paragraph standard-versus-policy note.
\end{enumerate}

\section{Lab 6: Initialization and aggregates}
\textbf{Standards map:} \refcl{6.7.11}.\\
\textbf{Learning goal:} Design robust initialization.\\
\textbf{Procedure:} Initialize nested structs using designated and empty initialization.\\
\textbf{Acceptance evidence:} Every field has expected state.
\begin{enumerate}
\item Create a clean branch and save the compiler/runtime version and invocation.
\item Build one success case and one intentionally failing test case.
\item Record expected diagnostic or runtime behavior without asserting unsupported guarantees.
\item Deliver source, commands, test output, and a one-paragraph standard-versus-policy note.
\end{enumerate}

\section{Lab 7: Type inference and constexpr}
\textbf{Standards map:} \refcl{6.7.10, 6.7.2}.\\
\textbf{Learning goal:} Understand C23 additions and constraints.\\
\textbf{Procedure:} Introduce auto, typeof and constexpr in isolated modules.\\
\textbf{Acceptance evidence:} Gated compile probes pass on supported toolchains.
\begin{enumerate}
\item Create a clean branch and save the compiler/runtime version and invocation.
\item Build one success case and one intentionally failing test case.
\item Record expected diagnostic or runtime behavior without asserting unsupported guarantees.
\item Deliver source, commands, test output, and a one-paragraph standard-versus-policy note.
\end{enumerate}

\section{Lab 8: Bit-precise and enum ABI}
\textbf{Standards map:} \refcl{6.7.3}.\\
\textbf{Learning goal:} Identify representation boundaries.\\
\textbf{Procedure:} Test \_BitInt and fixed underlying enum with serialization.\\
\textbf{Acceptance evidence:} Codec output identical across supported targets.
\begin{enumerate}
\item Create a clean branch and save the compiler/runtime version and invocation.
\item Build one success case and one intentionally failing test case.
\item Record expected diagnostic or runtime behavior without asserting unsupported guarantees.
\item Deliver source, commands, test output, and a one-paragraph standard-versus-policy note.
\end{enumerate}

\section{Lab 9: Attributes and diagnostics}
\textbf{Standards map:} \refcl{6.7.13}.\\
\textbf{Learning goal:} Separate annotations from language semantics.\\
\textbf{Procedure:} Gate nodiscard and maybe\_unused using \_\_has\_c\_attribute.\\
\textbf{Acceptance evidence:} Compare diagnostics without assuming universal support.
\begin{enumerate}
\item Create a clean branch and save the compiler/runtime version and invocation.
\item Build one success case and one intentionally failing test case.
\item Record expected diagnostic or runtime behavior without asserting unsupported guarantees.
\item Deliver source, commands, test output, and a one-paragraph standard-versus-policy note.
\end{enumerate}

\section{Lab 10: Checked arithmetic}
\textbf{Standards map:} \refcl{7.20}.\\
\textbf{Learning goal:} Avoid overflow in size calculations.\\
\textbf{Procedure:} Implement checked\_add/mul and fallback.\\
\textbf{Acceptance evidence:} Boundary tests preserve preconditions.
\begin{enumerate}
\item Create a clean branch and save the compiler/runtime version and invocation.
\item Build one success case and one intentionally failing test case.
\item Record expected diagnostic or runtime behavior without asserting unsupported guarantees.
\item Deliver source, commands, test output, and a one-paragraph standard-versus-policy note.
\end{enumerate}

\section{Lab 11: Resource inclusion}
\textbf{Standards map:} \refcl{6.10.4}.\\
\textbf{Learning goal:} Use \#embed responsibly.\\
\textbf{Procedure:} Embed short binary test fixture under a feature gate.\\
\textbf{Acceptance evidence:} Build hashes and byte counts match expectations.
\begin{enumerate}
\item Create a clean branch and save the compiler/runtime version and invocation.
\item Build one success case and one intentionally failing test case.
\item Record expected diagnostic or runtime behavior without asserting unsupported guarantees.
\item Deliver source, commands, test output, and a one-paragraph standard-versus-policy note.
\end{enumerate}

\section{Lab 12: Formatted I/O}
\textbf{Standards map:} \refcl{7.23}.\\
\textbf{Learning goal:} Protect against format/length errors.\\
\textbf{Procedure:} Create bounded formatter and input parser.\\
\textbf{Acceptance evidence:} Truncation and malformed input cause errors.
\begin{enumerate}
\item Create a clean branch and save the compiler/runtime version and invocation.
\item Build one success case and one intentionally failing test case.
\item Record expected diagnostic or runtime behavior without asserting unsupported guarantees.
\item Deliver source, commands, test output, and a one-paragraph standard-versus-policy note.
\end{enumerate}

\section{Lab 13: Allocator integration}
\textbf{Standards map:} \refcl{7.24.4}.\\
\textbf{Learning goal:} Manage ownership under failures.\\
\textbf{Procedure:} Build a growable buffer with simulated realloc failure.\\
\textbf{Acceptance evidence:} Original buffer survives failure and zero-size policy is explicit.
\begin{enumerate}
\item Create a clean branch and save the compiler/runtime version and invocation.
\item Build one success case and one intentionally failing test case.
\item Record expected diagnostic or runtime behavior without asserting unsupported guarantees.
\item Deliver source, commands, test output, and a one-paragraph standard-versus-policy note.
\end{enumerate}

\section{Lab 14: Unicode conversion}
\textbf{Standards map:} \refcl{6.2.9, 7.30}.\\
\textbf{Learning goal:} Handle encoding state.\\
\textbf{Procedure:} Convert UTF input with error handling and boundary tests.\\
\textbf{Acceptance evidence:} Invalid byte sequences cannot yield unvalidated text.
\begin{enumerate}
\item Create a clean branch and save the compiler/runtime version and invocation.
\item Build one success case and one intentionally failing test case.
\item Record expected diagnostic or runtime behavior without asserting unsupported guarantees.
\item Deliver source, commands, test output, and a one-paragraph standard-versus-policy note.
\end{enumerate}

\section{Lab 15: Atomics and threads}
\textbf{Standards map:} \refcl{5.1.2.4, 7.17, 7.28}.\\
\textbf{Learning goal:} Reason about happens-before.\\
\textbf{Procedure:} Implement producer/consumer with mutex or release/acquire.\\
\textbf{Acceptance evidence:} Thread sanitizer does not report data races in covered paths.
\begin{enumerate}
\item Create a clean branch and save the compiler/runtime version and invocation.
\item Build one success case and one intentionally failing test case.
\item Record expected diagnostic or runtime behavior without asserting unsupported guarantees.
\item Deliver source, commands, test output, and a one-paragraph standard-versus-policy note.
\end{enumerate}

\section{Lab 16: Target portability}
\textbf{Standards map:} \refcl{5.3.5, Annex J}.\\
\textbf{Learning goal:} Investigate implementation-defined behavior.\\
\textbf{Procedure:} Collect data-model, endianness, char-signedness and alignment evidence.\\
\textbf{Acceptance evidence:} Platform register captures relevant differences.
\begin{enumerate}
\item Create a clean branch and save the compiler/runtime version and invocation.
\item Build one success case and one intentionally failing test case.
\item Record expected diagnostic or runtime behavior without asserting unsupported guarantees.
\item Deliver source, commands, test output, and a one-paragraph standard-versus-policy note.
\end{enumerate}

\section{Lab 17: Secure code audit}
\textbf{Standards map:} \refcl{6.5, 7.24, 7.26}.\\
\textbf{Learning goal:} Combine tools with review.\\
\textbf{Procedure:} Review a supplied deliberately flawed parser.\\
\textbf{Acceptance evidence:} All injected high-risk defects have fixes and regression tests.
\begin{enumerate}
\item Create a clean branch and save the compiler/runtime version and invocation.
\item Build one success case and one intentionally failing test case.
\item Record expected diagnostic or runtime behavior without asserting unsupported guarantees.
\item Deliver source, commands, test output, and a one-paragraph standard-versus-policy note.
\end{enumerate}

\section{Lab 18: Capstone verification}
\textbf{Standards map:} \refcl{4, 5, 6, 7}.\\
\textbf{Learning goal:} Construct a reproducible evidence package.\\
\textbf{Procedure:} Deliver a binary record parser with safe allocation and serialization.\\
\textbf{Acceptance evidence:} Matrix tests, SAST and documented exceptions pass.
\begin{enumerate}
\item Create a clean branch and save the compiler/runtime version and invocation.
\item Build one success case and one intentionally failing test case.
\item Record expected diagnostic or runtime behavior without asserting unsupported guarantees.
\item Deliver source, commands, test output, and a one-paragraph standard-versus-policy note.
\end{enumerate}

\chapter{Capstone deliverable specification}
\section{Functional criteria}
Extend the decoder to produce a validated owned copy, support a maximum payload budget, and serialize a response using explicit byte-order conversions. Provide all of: reference parser, unit tests, negative tests, fuzz harness, CMake configuration, API documentation and a target environment report.
\section{Failure injection}
\begin{itemize}
\item Truncated header, maximum payload, overlong payload and an inconsistent declared size.
\item Failure at every allocation point (use a test allocator interface).
\item Empty input, null output, exact boundary lengths and repeated decoding.
\item Inconsistent encoding, if any text payload extension is implemented.
\item Build under at least two optimization settings, and two compiler families when practical.
\end{itemize}
\section{Definition of done}
\begin{enumerate}\item No known undefined behavior on test cases; assumptions documented separately.
\item Sanitizers pass supported paths; static analysis findings triaged with evidence.
\item Every rejected input leaves output in its documented state.
\item Every optional C23 feature has a compile-time capability check or minimum version requirement.
\item One reviewer can reproduce the outcome from a clean checkout.\end{enumerate}
\chapter{Portfolio evidence checklist}
Record module identifiers, source references, build logs, test traces, findings, remediation commits, manual-review signoff, supported targets, unresolved portability issues and rationale for any deviation from engineering policy.\par
\begin{warnbox}\textbf{Tool safety.} Deliberately invalid snippets should be compiled in dedicated negative-test targets and never executed outside controlled sanitizer experiments. Compiler optimization may remove or transform effects after undefined behavior; demonstrations are not predictive execution models.\end{warnbox}

\end{document}
